\chapter{Thảo luận}
\label{ch:discussion}

Chương này (1) diễn giải các kết quả thực nghiệm của
Chương~\ref{ch:experiments}, sau đó chưng cất kinh nghiệm kỹ thuật thu được
khi xây dựng và vận hành pipeline thành các bài học về (2) khám phá website
thật, (3) dựng chốt chặn cấu trúc cho LLM agent thay vì chỉ nhắc nhở qua
prompt, (4) induction các operator chất lượng cao, và (5) huấn luyện và
serving mô hình ngôn ngữ nhỏ, trước khi (6) tổng hợp thành các best
practice và (7) bàn về các nguy cơ đối với tính hợp lệ. Mỗi bài học theo
cùng một khuôn: một vấn đề gặp thật trong các run của chúng tôi, chẩn đoán,
cách khắc phục, và nguyên tắc tổng quát rút ra. Chi tiết đầy đủ các sự cố
được tập hợp ở Phụ lục~\ref{app:runs}.

\section{Diễn giải kết quả}
\label{sec:disc-results}

% TODO: làm mới mục này khi có kết quả ba website còn lại và các lần lặp
% thống kê (Mục 4.5-4.6).

Con số nổi bật nhất đến nay, cải thiện từ 5{,}56\% lên 6{,}67\% tỷ lệ
hoàn thành trên bộ task GitLab chính thức, cần được đọc thận trọng thay
vì đắc thắng. Về tuyệt đối, mức tăng là nhỏ, và với mỗi task chỉ chạy một
lần, nó phải được coi là tín hiệu định hướng cho tới khi các lần chạy lặp
của Mục~\ref{sec:exp-stats} hoàn tất. Điều khiến tín hiệu này đáng chú ý là
hồ sơ chi phí của nó: nó thu được từ vỏn vẹn 352 mẫu huấn luyện, thu thập
tự động từ một website không có nhãn, và bản thân việc fine-tune mất chưa
đầy nửa giờ trên một GPU. Nói cách khác, pipeline đã dịch chuyển một mô
hình nhỏ đúng hướng với lượng dữ liệu nhỏ đáng kinh ngạc; câu hỏi mở là xu
hướng này kéo dài tới đâu khi thu thập dữ liệu mở rộng quy mô, điều mà
kết quả liên-website ở Mục~\ref{sec:exp-main} sẽ bắt đầu trả lời.

Ba yếu tố nhiều khả năng đang chặn mức tăng hiện tại. Thứ nhất, cỡ dữ
liệu: 352 mẫu thấp hơn các kho fine-tune agent điển hình hai đến ba bậc,
nên mô hình chỉ thấy mỗi khuôn tương tác vài lần. Thứ hai, độ phủ: dữ liệu
huấn luyện thu từ một phần website, trong khi benchmark rút task từ toàn
site. Thứ ba, chính mô hình gốc: một mô hình 9B tham số fine-tune thuần
giám sát không thể được kỳ vọng khôi phục các năng lực (lập kế hoạch dài
hạn, tự sửa lỗi) mà trajectory chỉ minh họa thưa thớt. Không yếu tố nào
trong ba điều trên quy lỗi cho thiết kế pipeline; cả ba đều chỉ về khối
lượng và độ phủ dữ liệu như ràng buộc trói buộc, đúng ràng buộc mà
pipeline được sinh ra để nới lỏng.

Quan trọng không kém là những gì lần chạy đầu--cuối đã \emph{chứng minh}:
mọi công đoạn (khám phá, phát hiện, ghép chuỗi operator, tinh lọc, huấn
luyện, serving, đánh giá chính thức) đều chạy trên website thật không
cần can thiệp tay, và các cổng chất lượng đã chặn được những dữ liệu lẽ ra
sẽ làm bẩn huấn luyện (Mục~\ref{sec:disc-slm}). Phần còn lại của chương kể
lại cái giá để câu vừa rồi đúng.

\section{Bài học khi khám phá website thật}
\label{sec:disc-exploration}

\subsection{Frontier cạn kiệt}

\textbf{Vấn đề.} Một run khám phá sớm trên website kiểu Reddit khựng lại ở
57 trên ngân sách 500 node: frontier đơn giản là cạn.
\textbf{Chẩn đoán.} Khi đó frontier chỉ lớn lên từ các URL solver tình cờ
điều hướng tới. Trang nào mà mọi task đề xuất đều thuộc dạng tra cứu thì
không sinh điều hướng nào, nên nguyên vùng website không bao giờ vào hàng
đợi; khám phá bị bắt làm con tin bởi việc explorer tình cờ đề xuất gì.
\textbf{Khắc phục.} Bổ sung khai thác hyperlink làm tầng frontier thứ hai,
kèm URL template và trần mỗi-template để trang gần trùng không tràn hàng
đợi (Mục~\ref{sec:method-frontier}).
\textbf{Bài học.} Trên website thật, \emph{khả năng chạm tới không được
phụ thuộc hành vi giải task}. Tín hiệu cấu trúc (link) và tín hiệu hành vi
(điều hướng) là hai nguồn frontier bổ trợ, và nguồn cấu trúc cần kiểm soát
đa dạng tường minh vì nó đến theo lô.

\subsection{Đề xuất task bị suy sụp}

\textbf{Vấn đề.} Trong một run thực tế, 46\% tổng số task đề xuất gần như
trùng nhau giữa các trang; các biến thể của ``vào trang Diễn đàn'' được
đề xuất lại ở hầu như mọi node.
\textbf{Chẩn đoán.} Mọi trang của một website chung menu điều hướng toàn
cục, còn một explorer bị prompt theo từng trang thì không nhớ mình đã đề
xuất gì ở nơi khác.
\textbf{Khắc phục.} Explorer giờ tồn tại suốt run: các đề xuất gần đây được
đưa ngược vào prompt làm ngữ cảnh phủ định, các bản gần trùng bị loại thêm
ở mức mã (Mục~\ref{sec:method-lowlevel}).
\textbf{Bài học.} Sự đa dạng của task sinh ra là thuộc tính \emph{cấp
run}, không phải cấp trang; nó phải được thi hành bằng trạng thái sống lâu
hơn bất kỳ lời gọi LLM đơn lẻ nào.

\subsection{Trang không đáng tin}

\textbf{Vấn đề.} Một trang chậm duy nhất, timeout điều hướng trên một
URL khai thác được, làm sập nguyên một run khám phá nhiều giờ.
\textbf{Chẩn đoán.} Các thao tác trình duyệt được tin tưởng ngầm định; một
ngoại lệ không bắt ở bất kỳ đâu trong xử lý node lan thẳng lên đỉnh.
\textbf{Khắc phục.} Mọi thao tác trình duyệt giờ đều fail-soft: node không
mở được thì ghi log, lưu kèm lỗi, rồi bỏ qua; còn ngoại lệ thật sự bất ngờ
vẫn làm run dừng có tiếng động. Sản phẩm trung gian được lưu sau mỗi node,
nên run bị ngắt vẫn để lại bản chụp nhất quán, nối lại được.
\textbf{Bài học.} Hãy coi mọi tương tác với website thật là thao tác
\emph{rồi sẽ} thất bại, và quyết định cho từng thao tác xem thất bại có
sống sót được không; làm cho tiến độ của run bền vững ở mức hạt mịn nhất
ngân sách cho phép.

\subsection{Trạng thái dùng chung, có thể bị sửa}

\textbf{Vấn đề.} Website benchmark là tài nguyên dùng chung và có trạng
thái: một agent bất cẩn có thể làm bẩn (hoặc phá hủy) môi trường mà
các run khác phụ thuộc.
\textbf{Khắc phục.} Ba lớp kỷ luật. Lớp thứ nhất là chế độ read-only, loại
các task thay đổi trạng thái theo nguyên tắc fail-closed: task không phân
loại được thì bị coi là có thay đổi. Lớp thứ hai là một luật cứng cấm các
task phá hoại cấp tài khoản bất kể chế độ. Lớp thứ ba áp dụng khi đánh giá
trên instance sống dùng chung: task benchmark được phân loại theo ý định
thành hai tập read-only và read-write, để hai máy có thể quét cùng website đồng
thời mà không giẫm chân nhau (Mục~\ref{sec:disc-slm}).
\textbf{Bài học.} Trên hạ tầng dùng chung, an toàn phải
\emph{fail-closed}: gánh nặng chứng minh thuộc về việc một task vô hại,
chứ không phải việc nó có hại.

\section{Chốt chặn cấu trúc thay vì nhắc nhở qua prompt}
\label{sec:disc-guardrails}

Một chủ đề lặp lại trong các run của chúng tôi: bảo mô hình cư xử đúng
không đồng nghĩa hệ thống cư xử đúng. Các sự cố dưới đây đều khởi đầu là
một quy ước mức prompt bị mô hình thật lặng lẽ vi phạm, và đều kết thúc
bằng một cơ chế cấu trúc không phụ thuộc sự vâng lời của mô hình.

\subsection{Quy ước gắn thẻ không sống sót khi gặp mô hình thật}

\textbf{Vấn đề.} Đề xuất task ban đầu được phân loại bằng thẻ
\texttt{[INFO]}\slash\texttt{[NAV]}\slash\texttt{[MOD]} đặt đầu câu mà
explorer được chỉ thị phải ghi. Trong một run thực tế, 89\% đề xuất không mang thẻ nào.
\textbf{Chẩn đoán.} Xu hướng tự nhiên của mô hình là trả phân loại trong
một trường \texttt{type} có cấu trúc (thứ mà parser vứt bỏ) thay vì
một cặp ngoặc trong chuỗi. Tệ hơn, phép kiểm tra an toàn cho task thay đổi
trạng thái là so khớp chuỗi ký tự trên \texttt{[MOD]}, nên một task phá
hoại không gắn thẻ sẽ lọt qua êm ái.
\textbf{Khắc phục.} Đề xuất giờ được yêu cầu và phân tích như đối tượng có
cấu trúc với trường phân loại tường minh; thứ gì không phân loại được gắn
\texttt{UNKNOWN} và bị loại trong chế độ read-only.
\textbf{Bài học.} Hãy xin mô hình \emph{cấu trúc}, đừng xin quy ước định
dạng bên trong văn bản tự do, và tuyệt đối đừng xây phép kiểm tra an
toàn trên một quy ước mà mô hình có thể không theo.

\subsection{Vòng kẹt, qua hai lần lặp}

\textbf{Vấn đề.} Solver đốt trọn ngân sách bước để lặp một hành động rõ
ràng không có tác dụng; trong một episode ghi lại được, nó click cùng
một phần tử trơ suốt cả mười lăm bước.
\textbf{Chẩn đoán và khắc phục lần một.} Chúng tôi thêm một cảnh báo, tiêm
vào prompt sau hai hành động liên tiếp không hiệu ứng trên cùng đích. Run
thật kế tiếp bác bỏ cách sửa này: mô hình tiếp tục click dù cảnh báo hiện
diện từ lời gọi thứ ba trở đi. Nó cũng lộ ra một điểm tinh vi: mô hình
đôi khi thử lại một dropdown bị kẹt bằng các đối số trông-khác-nhưng-sai-
như-nhau, nên so trùng nguyên văn chuỗi hành động sẽ bỏ lọt vòng kẹt.
\textbf{Khắc phục lần hai.} Bộ phát hiện kẹt giờ khóa theo \emph{đích} bị
tác động và theo việc quan sát có đổi không; sau một số lần lặp không tiến
triển có giới hạn, episode bị chấm dứt thẳng, không tốn thêm lời gọi LLM
nào. Yêu cầu ``không thay đổi quan sát được'' giữ chốt chặn khỏi kích hoạt
nhầm trên các chuỗi hợp lệ tái dùng một phần tử qua các trang thay đổi.
\textbf{Bài học.} Cảnh báo trong prompt là một lời gợi ý; ngân sách là một
cơ chế. Khi mô hình chứng minh nó phớt lờ gợi ý, hãy leo thang thành cơ chế
và kiểm chứng bước leo thang trên hành vi thật, không phải trên hành vi
mà prompt yêu cầu.

\subsection{Các sự cố nhỏ hơn, cùng một khuôn}

Ba sự cố nữa theo đúng khuôn trên và nhận cùng kiểu sửa. Task tra cứu do
explorer đề xuất đôi khi chứa sẵn đáp án (``tìm bài đăng có tiêu đề X''),
thứ solver ``giải'' được mà không học gì; khâu sinh giờ có chỉ thị chống
điều đó và các task lộ đáp án bị lọc. Lời gọi có đính kèm ảnh thỉnh thoảng
trả về rỗng (từ chối phía thượng nguồn); agent thử lại một lần ở chế độ
thuần văn bản thay vì hỏng cả bước. Và các episode cạn ngân sách trước đây
kết thúc bằng một thất bại câm lặng. Giờ chúng nhận một lời gọi wrap-up
cuối, thuần văn bản, mà đầu ra hợp lệ \emph{duy nhất} là tuyên bố hoàn
thành hoặc báo cáo bất khả thi, biến một lần hết giờ im lặng thành một
kết cục có nhãn mà dataset dùng được.

\section{Bài học từ induction operator}
\label{sec:disc-operators}

Cơ chế semantic task của Mục~\ref{sec:method-semantic} chỉ tốt ngang các
operator nó ghép. Hai chế độ hỏng lộ ra khi chúng tôi kiểm toán các
operator sinh trong một run thực tế sớm, và cả hai đều dẫn tới các cổng
chặn cấu trúc thay vì chỉnh prompt.

\subsection{Operator tautology}

\textbf{Vấn đề.} Ba trong bốn operator sinh ra trong run đó có tập effect
trùng hệt tập precondition.
\textbf{Chẩn đoán.} Khi được yêu cầu mô tả precondition và effect của một
skill, LLM thấy dễ nhất là nhắc lại bối cảnh: ``đang ở danh sách issue''
trước, ``đang ở danh sách issue'' sau. Operator như vậy không mã hóa thay
đổi trạng thái nào, nên không bao giờ đứng giữa chuỗi và không dẫn tới đâu
mới; nếu bỏ mặc, tautology sẽ lấp đầy đồ thị toán tử bằng các node chết.
\textbf{Khắc phục.} Cổng chặn tautology: effect nào lặp lại precondition bị
bỏ, và skill không còn effect mới bị từ chối nâng cấp, có ghi log lý do.
\textbf{Bài học.} Với các sản phẩm ký hiệu do LLM viết, \emph{hãy định
nghĩa tính hợp lệ theo lý thuyết tập hợp và thi hành nó bằng mã}; đừng
trông cậy prompt cấm được các câu trả lời thoái hóa.

\subsection{Mệnh đề đồng nghĩa mà embedding không gộp nổi}

\textbf{Vấn đề.} Ghép chuỗi đòi hỏi effect của operator này khớp
precondition của operator kia, nhưng các operator induction độc lập diễn
đạt cùng một điều kiện theo những cách khác nhau.
\textbf{Chẩn đoán.} Chúng tôi đo độ tương đồng embedding của các cặp mệnh
đề trong cùng run: một cặp diễn đạt tương đương thật chấm 0{,}570, 
\emph{dưới} ngưỡng chuẩn hóa 0{,}6, trong khi một cặp khác nghĩa chấm
0{,}606, trên ngưỡng. Không ngưỡng nào tách được hai trường hợp; chuẩn hóa
thuần hình học không thể khôi phục một bộ từ vựng chung sau khi sự đã rồi.
\textbf{Khắc phục.} Phòng thay vì chữa: induction nhận biết từ vựng. Prompt
trình các nhãn mệnh đề sẵn có của registry kèm chỉ thị tái sử dụng nguyên
văn mệnh đề tương đương, để việc thống nhất diễn ra lúc \emph{viết} (nơi
mô hình có đủ ngữ cảnh) thay vì lúc so embedding, nơi chỉ còn bề mặt câu
chữ.
\textbf{Bài học.} Khi hai đầu ra LLM về sau phải \emph{khớp} nhau, hãy chia
sẻ bộ từ vựng ngay lúc sinh; đừng ủy thác việc thống nhất ngữ nghĩa cho một
thước đo tương đồng.

\section{Cạm bẫy khi huấn luyện và serving SLM}
\label{sec:disc-slm}

\subsection{Chat template là một phần của mô hình}

Họ mô hình đích mặc định phát các đoạn suy nghĩ tường minh, trong khi đích
huấn luyện của chúng tôi là các cặp thought-action trần. Chat template
đúng (đã tắt thinking) phải được tìm từ registry template của framework
huấn luyện thay vì đoán theo tài liệu, và cùng thiết lập đó phải được áp
lại, trên từng request, lúc serving. Lệch ở bất kỳ phía nào là âm thầm huấn
luyện hoặc đánh giá một phân phối khác với dự định. Điểm tổng quát:
\emph{template và cấu hình giải mã là một phần của sản phẩm mô hình}, và
cấu hình lúc huấn luyện với lúc serving phải được ghim cùng nhau.

\subsection{Lỗi prompt-fidelity}

Thất bại nhiều bài học nhất của khâu huấn luyện không sinh ra lỗi nào. Một
bài kiểm nhanh sau huấn luyện, viết bằng một prompt \emph{xấp xỉ} prompt
của solver, khiến mô hình fine-tune phát ra mô tả bằng lời thay vì hành
động thực thi được, trông y hệt một lần fine-tune hỏng. Mô hình không
sao cả; prompt kiểm thử mới là thứ không trung thực. Kiểm lại bằng prompt
sản xuất (import từ solver, không gõ lại) cho hành động đúng. Sự cố
này là lý do prompt parity (Mục~\ref{sec:method-data}) được thi hành bằng
cách \emph{import} system prompt thật và các hàm phụ trợ của nó vào trình
dựng dataset lẫn bộ đánh giá, không bao giờ bằng cách chép tay. Mọi bài
kiểm một agent fine-tune chỉ đáng tin ngang độ trung thực của prompt kiểm
so với sản xuất.

\subsection{Cách ly phụ thuộc}

Hỗ trợ một kiến trúc mô hình mới phát hành đòi hỏi bản phát triển của thư
viện mô hình hóa, xung đột với yêu cầu ghim phiên bản của engine serving.
Ba môi trường cách ly (huấn luyện, serving, đánh giá) giải quyết điều
mà không môi trường đơn nào thỏa mãn nổi. Tầm thường, nhưng nó biến một
xung đột phiên bản vô phương thành một bước cài đặt thường quy.

\subsection{Agent tự chấm bài mình}

Bước judge được đưa vào vì thành công do solver tự tuyên bố không thể được xem là bằng chứng đáng tin cậy về việc hoàn thành tác vụ. Khi solver trả lời rằng tác vụ đã “hoàn tất”, nó đồng thời tự đánh giá kết quả của chính mình. Vì vậy, các episode được tuyên bố thành công cần được đánh giá lại dựa trên trạng thái cuối thực tế.

Cơ chế này giúp phát hiện những trường hợp solver tự kết luận sai rằng tác vụ đã thành công trước khi các trajectory tương ứng được đưa vào tập dữ liệu huấn luyện. Trong chế độ nghiêm ngặt, trình dựng dataset mặc định loại bỏ mọi trajectory được tuyên bố thành công nhưng chưa qua judge, thay vì mặc nhiên chấp nhận chúng.

Theo nguyên tắc P1, \emph{dữ liệu chưa được kiểm chứng không phải là dữ liệu trung tính; đó là một khoản nợ chất lượng có thể bị khuếch đại trong quá trình huấn luyện}.

\subsection{Đánh giá an toàn trên website sống dùng chung}

Chạy hai lượt đánh giá (gốc và fine-tune) trên một instance website sống
bình thường sẽ phải xếp hàng trọn một ngày làm việc. Phân loại task
benchmark theo ý định thành hai tập read-only và read-write cho phép máy thứ hai
đánh giá tập read-only song song với lượt quét không hạn chế của máy thứ
nhất. Bản thân việc phân loại chứa hai cái bẫy đáng ghi lại: va chạm cụm
danh từ (``merge request'' chứa động từ chỉ về mặt cú pháp; ``report bugs''
không gọi tên hành động báo cáo mà regex chờ đợi) và việc chèn từ đánh bại
các mẫu so cụm ngây thơ (``open \emph{a new} issue''). Chi phí truyền tải
cũng đáng kể: chuyển adapter LoRA 1{,}7\,GB rồi gộp lại tại chỗ thắng việc
chuyển checkpoint đã gộp 18\,GB.

\section{Best practice}
\label{sec:disc-practices}

Bảng~\ref{tab:practices} cô đọng các bài học của chương thành những thực
hành chúng tôi sẽ áp dụng từ ngày đầu của một dự án tương tự.

\begin{table}[H]
\centering
\caption[Các best practice đã chưng cất]{Các best practice cho pipeline
thu thập dữ liệu bằng LLM agent trên website thật.}
\label{tab:practices}
\small
\begin{tabular}{@{}L{4.4cm} L{8.8cm}@{}}
\toprule
\textbf{Thực hành} & \textbf{Lý do (mục)} \\
\midrule
Kiểm chứng rồi mới thu nhận & Skill, operator, thành công đều là khoản nợ
cho tới khi được kiểm độc lập (\ref{sec:disc-slm}, \ref{sec:disc-operators}) \\
Prompt parity bằng import & Dữ liệu huấn luyện, agent thật, bộ đánh giá
dùng chung một nguồn prompt; chép tay sẽ trôi dạt (\ref{sec:disc-slm}) \\
Cấu trúc thay vì quy ước & Xin mô hình trường có kiểu; đừng bao giờ trích
thông tin an toàn từ văn bản tự do (\ref{sec:disc-guardrails}) \\
Cơ chế thay vì cảnh báo & Chuyển vi phạm lặp lại của mô hình thành ngân
sách và lệnh dừng cứng, kiểm chứng trên run thật (\ref{sec:disc-guardrails}) \\
Duyệt web fail-soft & Thao tác trang nào cũng có thể hỏng; bỏ qua và ghi
log cục bộ, chỉ sập với điều thật sự bất ngờ (\ref{sec:disc-exploration}) \\
Lưu bền theo từng node & Run bị ngắt phải để lại sản phẩm nhất quán, nối
lại được (\ref{sec:disc-exploration}) \\
An toàn fail-closed & Task không phân loại được tính là phá hoại; hành động
phá hoại cấp tài khoản bị cấm thẳng (\ref{sec:disc-exploration}) \\
Chia sẻ từ vựng lúc sinh & Sản phẩm ký hiệu phải khớp nhau về sau nên được
sinh trên một bộ nhãn chung (\ref{sec:disc-operators}) \\
Ghim template cùng mô hình & Chat template và cờ giải mã là một phần của
sản phẩm, cho cả huấn luyện lẫn serving (\ref{sec:disc-slm}) \\
Ghi mọi thứ thành sự kiện & Nhật ký sự kiện từng run là thứ khiến phân
tích hậu kỳ (phần lớn chương này) khả thi (\ref{sec:disc-exploration}) \\
\bottomrule
\end{tabular}
\end{table}

\section{Nguy cơ đối với tính hợp lệ}
\label{sec:disc-validity}

\textbf{Hợp lệ nội tại.} Đánh giá đến nay chạy mỗi task benchmark một lần
mỗi mô hình; kết cục từng task trên website sống có nhiễu, nên khoảng cách
đo được có thể dịch chuyển khi lặp lại. Các lần chạy lặp dự kiến
(Mục~\ref{sec:exp-stats}) xử lý trực diện điều này. Vài cổng chất lượng
(judge, bộ viết lại suy nghĩ) bản thân là LLM và thừa kế thiên kiến LLM;
chúng tôi giảm thiểu nhưng không loại trừ bằng cách đối chiếu phán quyết
với trạng thái cuối cụ thể thay vì chỉ với bản ghi hội thoại.

\textbf{Hợp lệ ngoại vi.} Toàn bộ dữ liệu và đánh giá đến từ các website
WebArena tự triển khai: thực tế, nhưng là bốn website của một họ
benchmark, và từ một mô hình 9B duy nhất. Các tuyên bố được giới hạn
tương ứng: chúng tôi báo cáo pipeline cải thiện \emph{mô hình này} trên
\emph{các website này}, và coi khái quát hóa rộng hơn là việc tương lai
thay vì một hệ quả ngầm định.

\textbf{Hợp lệ cấu trúc.} Tỷ lệ hoàn thành tác vụ theo bộ chấm chính thức là thước đo tiêu chuẩn trong lĩnh vực, nhưng đây là một chỉ số nhị phân và chỉ phản ánh trạng thái cuối cùng. Do đó, chỉ số này không ghi nhận mức độ tiến bộ từng phần trong quá trình thực hiện. Bên cạnh đó, các semantic task được xác nhận thông qua tự báo cáo của solver kết hợp với judge ngoại tuyến có mức độ tin cậy thấp hơn so với các tác vụ được kiểm chứng trực tiếp bằng chương trình. Vì vậy, chúng tôi lưu tách biệt hai chế độ kiểm chứng trong dữ liệu, qua đó cho phép thực hiện các phân tích nghiêm ngặt hơn về sau.

\vspace{1em}
\noindent
Các bài học này, cùng những kết quả mà chúng đặt trong ngữ cảnh, khép lại
phần thực nghiệm của khóa luận. Chương cuối tổng kết những gì đã xây dựng
và học được, và nơi công việc này sẽ đi tiếp.
